\documentclass[11pt]{article}
\emergencystretch=2em
\usepackage[margin=1in]{geometry}
\usepackage[T1]{fontenc}
\usepackage{lmodern,amsmath,amssymb,amsthm,booktabs,array,enumitem}
\usepackage[colorlinks=true,linkcolor=blue,citecolor=blue,urlcolor=blue]{hyperref}
\newtheorem{theorem}{Theorem}
\newtheorem{proposition}[theorem]{Proposition}
\newtheorem{corollary}[theorem]{Corollary}
\newtheorem{lemma}[theorem]{Lemma}
\theoremstyle{definition}
\newtheorem{definition}[theorem]{Definition}
\newtheorem{example}[theorem]{Example}
\theoremstyle{remark}
\newtheorem{remark}[theorem]{Remark}
\newcommand{\Prb}{\mathbb P}
\newcommand{\E}{\mathbb E}
\newcommand{\ind}{\mathbf 1}
\newcommand{\Bin}{\mathrm{Bin}}
\newcommand{\cx}{\le_{\mathrm{cx}}}
\newcommand{\st}{\le_{\mathrm{st}}}
\newcommand{\TV}{d_{\mathrm{TV}}}
\title{Risk-Gated Mechanization:\\Severity, Verifiability, and Decisions Invariant\\to the Executor's Error Law}
\author{Ivo Matija\v{s}evi\'c}
\date{October 2026}
\newcommand\blfootnote[1]{\begingroup\renewcommand\thefootnote{}\footnote{#1}\addtocounter{footnote}{-1}\endgroup}
\begin{document}
\maketitle
\blfootnote{\textcopyright~2026 Ivo Matija\v{s}evi\'c.\ This work is licensed under a Creative Commons
Attribution 4.0 International (CC BY 4.0) license, \url{https://creativecommons.org/licenses/by/4.0/}.\ Preprint, DOI \href{https://doi.org/10.5281/zenodo.23324667}{10.5281/zenodo.23324667}.}
\begin{abstract}
When should a step of agentic work be replaced by a deterministic mechanism? Earlier papers in
this series answered with escape probabilities and costs, and treated every escaped fault as
equally bad. We add severity and formalize a two-way principle, which reads informally as:
mechanize a high-severity step whose correctness can be verified, and do not mechanize one whose
correctness cannot. Read literally, the principle overreaches; we give its exact conditions.
Severity sets a risk budget; it does not decide alone. An option is admissible when its
worst-case loss over the models compatible with the evidence fits the budget. Mechanization is
mandatory when closure is admissible and every other available option is not, and forbidden when
a compatible closed model exceeds the budget; when no option is admissible, the step needs more
evidence, an exogenous gate, or a redesign. Verification is a sufficient way to make closure
admissible, and an evidence-compatible closed model whose expected loss exceeds the budget, such as an
indistinguishable unsafe twin, forbids it. The bounds that decide are an executor's conditional escape floor and ceiling and
a mechanism's certified ceiling. They hold uniformly over every error law, dependence
structure, adaptive routing, and mixture of executors that satisfies them; a stochastic-dominance
theorem extends this invariance to every nondecreasing loss. When evidence-indistinguishable worlds
prefer opposite actions, every evidence-based rule has worst-case regret that grows linearly with
severity. A martingale budget gives a scheduler an
anytime guarantee under adaptive routing. Persistence enters through convex order: within
conditionally independent processes, an uncertain mechanism fault mass is worse than fresh
errors of the same mean for every convex loss, by the convexity of Bernstein polynomials, while
the probability of any escape moves the other way. Production mixing averages floors;
conjunctive review intersects blind sets, and one reviewer dominates another for every producer
exactly when its blind set is contained in the other's. With linear loss and constant severity,
the results reduce to the earlier mechanization criterion. The decision-theoretic and probabilistic
tools are classical; the contribution is their application to the mechanization decision. These are
mathematical models, not measurements.
\end{abstract}

\section{Introduction}
Part~I of this series models gated computation and shows that a gate family with a common blind
set leaves every executor with a verification floor~\cite{p1}. Part~II shows that gates generated
within one family inherit its blind set, and that exogenous or exact evidence is needed to lower
the floor~\cite{p2}. Part~III treats escape, cost, and correlated escape~\cite{p3}. Part~IV studies
progressive semantic mechanization: moving work from executor calls to deterministic
mechanisms~\cite{p4}. It shows that closure helps only when the mechanism's escape-plus-fallback
risk is below the open path's risk, and that evidence cannot separate \emph{twins}, mechanisms
with the same observation law. Part~V removes the constant-fault-rate assumption and works with
conditional risks~\cite{p5}.

None of these papers weighs how bad a fault is. This paper adds severity through the following
principle.

\begin{quote}
\textbf{Risk-gated mechanization.} Consequence-weighted risk is the main gate for mechanization.
If a step's correctness is verifiable and its severity is high, the step must be mechanized. If
its correctness is not verifiable and its severity is high, it must not be mechanized; it stays
open behind a human or exact-evidence gate. A fully determined step is the simplest verifiable
case.
\end{quote}
Read literally, the principle is a heuristic. Verifiable closure is not mandatory if another option,
such as an exact added gate, also fits the budget; and failure to verify does not forbid closure if
every compatible closed model is safe. Theorem~\ref{thm:admissible} gives the exact conditions.

A second aim shapes the results. Executors differ: a person, and several language models at
different effort levels, have different and unknown error laws, and work is routed among them.
A scheduler's decision rule should not depend on knowing those laws. We therefore look for
statements that use only bounds, and that hold uniformly over every error law satisfying them.
\emph{Invariance} in this paper always means this uniformity. Without any bound, no positive floor
and no mandate to mechanize can exist, because a fault-free executor is among the possible laws.

\paragraph{Contributions.}
\begin{enumerate}[leftmargin=*,itemsep=2pt]
\item Classes of executor processes defined by full-history conditional escape bounds, closed
under mixtures and adaptive routing, with the classical binomial domination of such
processes~\cite{doerr} (Lemma~\ref{lem:mixture}, Theorem~\ref{thm:dominance}).
\item Robust admissibility under a severity-derived budget, which gives ``must'' and ``must not''
a precise meaning, including the case where no option is admissible (Theorem~\ref{thm:admissible}).
\item The regret price of the second arm: for observational twins, Le~Cam's two-point bound and
minimax-regret allocation give a worst-case regret linear in severity (Theorem~\ref{thm:regret}).
\item Expected-loss thresholds for both arms, with the actual-risk, ex-ante, and worst-case readings
kept apart, and the classical sharp convex bounds from the mean applied to mean-only certificates
(Theorem~\ref{thm:arms}, Proposition~\ref{prop:meanonly}).
\item Persistence versus fresh errors as a convex-order comparison of mixed binomial counts
(Theorem~\ref{thm:convex}).
\item An anytime consequence budget for an adaptive scheduler, from the standard supermartingale
construction (Theorem~\ref{thm:budget}).
\item Production mixing, reviewer conjunction, and the statewise-dominance condition for universal
reviewer superiority, in the gated form of design-diversity models (Proposition~\ref{prop:review}).
\item Structural severity as a weighted reach, with a damage envelope and an ancestry bound
(Proposition~\ref{prop:severity}).
\item A reduction to Part~IV's criterion (Proposition~\ref{prop:reduction}).
\end{enumerate}
The probabilistic and decision-theoretic tools are classical: coupling and domination by independent
variables, stochastic and convex orders, Bernstein polynomials, total variation and two-point bounds,
minimax risk and minimax regret, nonnegative supermartingales, and positive association. The
contribution is their application to the mechanization decision and the invariance they give.
Section~\ref{sec:related} gives the sources.

\section{Setting}\label{sec:setting}
\paragraph{Units, escapes, and losses.}
A step of work is executed on $n$ units, where $n$ is fixed unless stated otherwise. Let
$J_i\in\{0,1\}$ indicate that unit $i$ emits an escaped fault, and $N_n=\sum_{i\le n}J_i$. Let
$\mathcal F_{i-1}$ be the full history before unit $i$: earlier outputs and escapes, verdicts,
routing choices, and internal states. A \emph{loss function} is a nondecreasing
$\varphi:\{0,\dots,n\}\to[0,\infty)$ with $\varphi(0)=0$. Examples are the expected count
$\varphi(k)=k$, the any-escape indicator $\varphi(k)=\ind\{k\ge1\}$ of Part~V's success metric,
and convex losses that penalize concentrated damage. A step has severity $S\ge0$
(Section~\ref{sec:severity}); its loss is $S\varphi(N_n)$.

\paragraph{Executor classes.}
\begin{definition}[Conditional escape class]\label{def:class}
For $0\le\ell\le u\le1$, the class $\mathcal P(\ell,u)$ contains every escape process with
\[
\ell\le\Prb(J_i=1\mid\mathcal F_{i-1})\le u\quad\text{almost surely, for every }i\le n .
\]
\end{definition}
The lower bound is a \emph{floor}. Under Part~I's model, with first-attempt badness $p(h)$ and a
stealth share $\sigma(h)$ that passes every gate at once, the floor at history $h$ is at least
$p(h)\sigma(h)$~\cite{p1,p5}. The upper bound is a \emph{ceiling}; independent majority gating
drives it down, but not below the floor~\cite{p1,p5}. The class says nothing about arrival laws,
dependence, retries, or the shape of the error law. The bounds must hold on the full history,
including histories after earlier escapes; a floor on correct prefixes only is weaker
(Example~\ref{ex:prefix}). Under routing, an executor's bound must hold on the histories and tasks
actually routed to it; an unconditional benchmark does not supply it.

\paragraph{Mechanisms and evidence.}
A mechanism $M$ maps inputs to outputs. Its fault mass $\pi(M)$ is the probability that an input
from the target law is mapped to a faulty output. With a retained gate and a fallback path, its
per-unit escape is $d(M)=e(M)+b(M)f$, as in Parts~IV and~V. We write $a(M)$ and $A(M)$ for lower and
upper bounds on its conditional per-unit escape that hold at every history. When deployment is
ungated, or the retained gate accepts every faulty output, inputs are independent draws from the target
law, and there is no fallback, both equal $\pi(M)$; repeated or adaptively selected inputs need not
satisfy this. A retained gate can make escape smaller than fault mass: an exact fail-closed gate on an
always-faulty mechanism gives escape $0$. With a gate or fallback, the escape bounds come from the gate
and fallback model. A gate that only suppresses outputs keeps escape at most the fault mass, so a
fault-mass bound remains an escape ceiling; a fallback can add escapes, and then validation of raw fault
mass does not by itself certify the closed path's escape.

A promotion rule sees evidence: the full record of an evidence experiment, including any adaptive
or sequential choices, and auxiliary randomness. Two mechanisms are \emph{twins} for that experiment if
they induce the same law for its full record. Twin status is relative to the experiment: equality of
one-sample record laws does not imply equality for two samples (two independent fair bits and two
identical fair bits agree on one bit and differ on two). Independent repetitions of an identical
per-sample record law preserve twin status at every sample size. Let $\mathcal C$ be the mechanisms
compatible with the stated assumptions. Two twin floors are needed, one for raw fault mass and one
for closed-path escape:
\[
\tau_\pi(M)=\sup\{\pi(M'):M'\in\mathcal C\text{ a twin of }M\},\qquad
\tau(M)=\sup\{a(M'):M'\in\mathcal C\text{ a twin of }M\}.
\]
Part~IV shows that a valid upper bound on fault mass, computed by a fixed rule from the record, is at
least $\tau_\pi(M)$ with high probability~\cite{p4,p5}; the same argument shows that a valid upper bound
on closed-path escape is at least $\tau(M)$ with high probability. Each certificate must be compared with
its own floor: a fault-free mechanism whose retained gate always rejects and whose fallback always emits
a fault has $\pi(M)=0$ but $a(M)=1$. Certifying escape needs validation that covers the whole closed
path, including the gate and fallback. Twin floors limit what evidence can certify; they are not lower
bounds on actual risk: a fault-free mechanism can have a faulty twin.

\begin{definition}[Verifiable]\label{def:verifiable}
Fix a level $U$ in advance. A step is \emph{verifiable at level $U$ with confidence $1-\eta$} if an
available build-and-evidence experiment promotes a mechanism with $\Prb(A(M)>U\mid
\text{promoted})\le\eta$.
\end{definition}
The level is fixed before the experiment; a data-dependent bound is not treated as having posterior
error $\eta$. The guarantee is about the randomized build-and-evidence experiment, averaged over
promoted builds; it is not a guarantee for every build compatible with the evidence, even with
$\eta=0$. Robust admissibility of a deployed build needs the uniform premise of
Theorem~\ref{thm:admissible}(iii). A fully determined step, whose decision tree is known, closed, and checked exactly on
every branch, is verifiable at level $U=0$ with $\eta=0$. Zero-failure validation on $m$ independent
target samples bounds the joint probability of promotion with fault mass above $U$ by $(1-U)^m$, the
classical zero-failure demonstration bound~\cite{littlewood93,butler93}; the conditional level divides
this by the promotion probability~\cite{p5}.

\paragraph{Gates and blind sets.}
A producer $e$ emits a bad unit with probability $p_e$; given badness, its fault has law $\mu_e$. A
gate accepts every fault in its blind set. A blind set is \emph{complete} if it is exactly the set of
accepted faults, and \emph{guaranteed} if it is only known to be accepted. A conjunctive stack
accepts a fault only if every member accepts it~\cite{p2}.

\section{Invariance lemmas}\label{sec:invariance}
\begin{lemma}[Mixtures and routing stay in the class]\label{lem:mixture}
Let $K$ be a component index, possibly chosen adaptively from the history, and suppose
$\ell\le\Prb(J_i=1\mid\mathcal F_{i-1},K)\le u$ almost surely for every $i$. Then the mixed process is
in $\mathcal P(\ell,u)$. In particular, adaptive routing among executors whose conditional escape
lies in $[\ell_e,u_e]$ on the routed histories yields a process in
$\mathcal P(\min_e\ell_e,\max_eu_e)$; with routing weights $w_{i,e}$, the conditional escape is
$\sum_ew_{i,e}q_{i,e}\ge\sum_ew_{i,e}\ell_e$.
\end{lemma}
\begin{proof}
By the tower property, $\Prb(J_i=1\mid\mathcal F_{i-1})=\E[\Prb(J_i=1\mid\mathcal F_{i-1},K)\mid
\mathcal F_{i-1}]$, an average of values in $[\ell,u]$.
\end{proof}

\begin{theorem}[Binomial dominance]\label{thm:dominance}
For every process in $\mathcal P(\ell,u)$,
\[
\Bin(n,\ell)\ \st\ N_n\ \st\ \Bin(n,u).
\]
Hence $\E\varphi(\Bin(n,\ell))\le\E\varphi(N_n)\le\E\varphi(\Bin(n,u))$ for every loss function
$\varphi$. For the any-escape indicator, a floor on correct prefixes alone gives Part~V's survival
bound $\Prb(N_n=0)\le(1-\ell)^n$.
\end{theorem}
\begin{proof}
This is the classical domination of a binary sequence with bounded conditional probabilities by
independent variables~\cite[Lemma~1.8.7]{doerr}, stated there for conditioning on past indicators,
which the tower property gives from bounds given $\mathcal F_{i-1}$. We include the coupling. On an
enlarged probability space, take independent uniform variables $V_1,\dots,V_n$ independent of
everything else. Construct the process sequentially: given the history, set
$J_i=\ind\{V_i\le\Prb(J_i=1\mid\mathcal F_{i-1})\}$ and draw the remaining history variables from
their conditional law given $J_i$. This reproduces the joint law. Put $L_i=\ind\{V_i\le\ell\}$ and
$H_i=\ind\{V_i\le u\}$. Then $L_i\le J_i\le H_i$ pointwise, with $\sum L_i\sim\Bin(n,\ell)$ and
$\sum H_i\sim\Bin(n,u)$. Pointwise order gives stochastic order, and $\varphi$ is nondecreasing. The
survival bound is Part~V's conditional floor.
\end{proof}

\begin{example}[A correct-prefix floor does not order counts]\label{ex:prefix}
Let $J_1$ be Bernoulli$(1/2)$ and $J_2=1-J_1$. The conditional risks on correct prefixes are $1/2$ and
$1$, both at least $1/2$, but $N_2=1$ always, which is not stochastically above $\Bin(2,1/2)$: the
binomial gives $\Prb(N\ge2)=1/4$. Likewise a marginal ceiling does not give a binomial ceiling:
$J_1=J_2=Z$ with $Z$ Bernoulli$(0.1)$ has $\Prb(N_2=2)=0.1>0.01$.
\end{example}

\begin{remark}
Any statement that uses only $\E\varphi(\Bin(n,\ell))$ and $\E\varphi(\Bin(n,u))$ therefore holds
uniformly over arrival laws, dependence, retry laws, and every mixture of executors within the
bounds. The map $t\mapsto\E\varphi(\Bin(n,t))$ is a nondecreasing polynomial, hence continuous. An
outcome-dependent number of units needs the stopping-time form of Section~\ref{sec:budget}.
\end{remark}

\section{Robust admissibility under a risk budget}\label{sec:admissible}
Severity enters as a budget. Fix a tolerated expected loss $K>0$ for a step of severity $S>0$. For an
option $a$, let $\mathcal C_a$ be the set of models compatible with the available evidence and
assumptions, and define its worst-case expected loss
\[
r_a^+=\sup_{P\in\mathcal C_a}\E_P\varphi(N_a).
\]
This is the maximum risk of Wald's minimax criterion~\cite{wald45}, taken over the set of compatible models, as in multiple-prior
decision theory and partial identification~\cite{gilboa89,manski00}. Option $a$ is \emph{robustly
admissible}\footnote{Here \emph{admissible} means that the option fits the budget. In statistical
decision theory the word means that a rule is not dominated~\cite{wald45}; the two notions are
unrelated.} if $Sr_a^+\le K$. Under linear loss, $r_a^+/n$ is the
worst-case per-unit escape, and admissibility reads $r_a^+/n\le r^\ast$ with the \emph{risk budget}
$r^\ast=K/(Sn)$. High severity means a small budget. For closure we distinguish two readings. A
\emph{deployed build} is robustly admissible if the worst case over builds compatible with the
evidence fits the budget. A \emph{build-and-evidence experiment} is admissible ex ante if the worst
case, over experiments satisfying Definition~\ref{def:verifiable}, of the expected loss averaged over
promoted builds fits the budget.

\begin{theorem}[Admissibility; the two arms]\label{thm:admissible}
Consider the options \emph{open} (executors with a gate stack), \emph{gated} (open with an added gate
$H$ from an available set), and \emph{closed} (a promoted mechanism).
\begin{enumerate}[label=(\roman*),itemsep=2pt]
\item \emph{Invariance.} Replacing $\mathcal C_a$ by its convex hull does not change $r_a^+$, so
whole-process mixtures of compatible models leave admissibility unchanged. If $\mathcal C_a$ is a bound
class $\mathcal P(\ell,u)$, it is also closed under adaptive routing (Lemma~\ref{lem:mixture}). A compatible
set given by specific executor laws need not be: routing can produce a process outside its convex hull.
\item \emph{Open path.} If the open path's compatible set is $\mathcal P(\ell,u)$, then
$r_{\mathrm{open}}^+$ lies between $\E\varphi(\Bin(n,\ell))$ and $\E\varphi(\Bin(n,u))$. If
$S\E\varphi(\Bin(n,\ell))>K$, the open path is inadmissible for every error law, dependence structure,
and routing within the bounds. A gate generated within the producing family leaves the floor unchanged
under Part~II's inheritance assumptions.
\item \emph{Closure, exact certificate.} If every build compatible with the evidence has its closed units
in $\mathcal P(0,U)$, the deployed build is robustly admissible whenever $S\E\varphi(\Bin(n,U))\le K$.
A fully determined step checked exactly on every branch gives this with $U=0$.
\item \emph{Closure, statistical certificate.} If the step is verifiable at level $U$ with confidence
$1-\eta$, and given promotion the closed units are in $\mathcal P(0,A(M))$, then the expected loss of a
promoted build, averaged over the experiment, is at most $(1-\eta)\E\varphi(\Bin(n,U))+\eta\varphi(n)$.
This gives ex-ante admissibility of the experiment, not robust admissibility of the deployed build: an
always-promoted build that is fault-free with probability $0.99$ and always faulty otherwise has
$U=0$ and $\eta=0.01$, average risk $0.01$, and worst-case risk $1$.
\item \emph{Closure forbidden.} If some compatible twin $M'$ of $M$ has closed units in
$\mathcal P(a(M'),1)$ and $S\E\varphi(\Bin(n,a(M')))>K$, closure is not robustly admissible. In
particular this holds when $S\E\varphi(\Bin(n,\tau(M)))>K$. If the twin relation holds for the whole
evidence experiment at every sample size, no amount of that evidence changes the conclusion.
\item \emph{Gated open path.} With complete blind sets, adding a gate $H$ gives producer $e$ the
first-attempt floor $p_e\mu_e(\Sigma\cap\Sigma_H)$. If $\Sigma_H$ is only guaranteed, this is a lower bound;
lowering the ceiling needs knowledge of what $H$ rejects. Exact evidence on a scope removes escape there
if it is sound, bound to the executed artifact or action, and enforced fail-closed. A human gate has no
special status: its effect is its acceptance set.
\end{enumerate}
Consequently, closure is \emph{mandatory} among the available options when it is admissible in the
chosen reading and the open path and every available gated option are inadmissible. Closure is
\emph{forbidden} when (v) holds. If no option is admissible, the step needs more distinguishing evidence,
an additional exogenous gate, a lower severity through redesign, or must not be performed.
\end{theorem}
\begin{proof}
(i) A supremum of expectations over a set equals the supremum over its convex hull, because an
expectation is linear in the model. Routing closure of $\mathcal P(\ell,u)$ is Lemma~\ref{lem:mixture}.
For the last claim, let executor A always produce escapes $(1,0)$ on two units and executor B always
$(0,1)$: every whole-process mixture has count $1$, while routing unit~1 to A and unit~2 to B gives $2$.
(ii) Theorem~\ref{thm:dominance} gives both bounds, and the lower bound applies to every compatible law.
A generated gate keeps the family's blind set (Part~II), so the first-attempt stealth share and the floor
are unchanged.
(iii) By Theorem~\ref{thm:dominance}, every compatible build has expected loss at most
$\E\varphi(\Bin(n,U))$.
(iv) On the event $A(M)\le U$, of conditional probability at least $1-\eta$, the closed count is
stochastically below $\Bin(n,U)$; otherwise the loss is at most $\varphi(n)$. The example has the stated
values.
(v) The twin is compatible and its expected loss is at least $\E\varphi(\Bin(n,a(M')))$ by
Theorem~\ref{thm:dominance}. The supremum over compatible twins of that quantity is
$\E\varphi(\Bin(n,\tau(M)))$, by monotonicity and continuity. If twins stay twins at every sample size,
the compatible set never excludes $M'$.
(vi) A fault passes the conjunction only if every member accepts it; with complete sets the passing set
is exactly $\Sigma\cap\Sigma_H$, and with guaranteed sets it contains it. Sound, bound, fail-closed exact
evidence accepts no fault on its scope.
\end{proof}

\begin{corollary}[Below the floor, inherited gates are not enough]\label{cor:exogenous}
If $S\E\varphi(\Bin(n,\ell_e))>K$ for every executor $e$, no open path whose gates all come from the
producing family is admissible, under any routing. An admissible option, if one exists, is a closure
with $Sr_{\mathrm{closed}}^+\le K$, or an open path gated outside the family. Parts (iii), (iv), and (vi)
give sufficient ways to establish these; they are not necessary. For example, with $n=2$, linear loss,
$S=1$, $K=0.05$, and an open escape of $0.1$ per unit, a closed process with $J_1=J_2=Z$ and
$Z\sim$ Bernoulli$(0.01)$ is robustly admissible ($\E N_2=0.02$), although its uniform conditional ceiling is
$1$, so neither (iii) nor (iv) certifies it.
\end{corollary}
\begin{proof}
By Theorem~\ref{thm:admissible}(ii) and Lemma~\ref{lem:mixture}, every routing of an open path with
inherited gates has expected loss at least $S\E\varphi(\Bin(n,\min_e\ell_e))>K$. The example has
$\E N_2=2\cdot0.01$ and $\Prb(J_2=1\mid J_1=1)=1$.
\end{proof}

\begin{remark}
The theorem gives the principle precise sufficient conditions, not unconditional implications.
\emph{Verifiable at a level within budget} is a sufficient way to make closure admissible; with a
statistical certificate the admissibility is ex ante. \emph{High severity} means a budget below the executor floors.
A step that is not verifiable at a chosen level is not thereby forbidden: closure is admissible when
every compatible closed model fits the budget, whether or not twins are separated. A safe twin set is
not enough, because compatible models need not be twins: a safe $M_0$ with evidence Bernoulli$(1/2)$ and
an unsafe $M_1$ with evidence Bernoulli$(1/4)$ both remain compatible after observing $0$. Closure is
forbidden when some compatible closed model violates the budget; a violating twin is one way this
happens. Severity alone decides nothing: it sets the budget, and the comparison
with the floors, the certified level, and the compatible twins decides. Part (v) concerns what evidence
permits, not actual risk; Theorem~\ref{thm:regret} prices the gap.
\end{remark}

\section{Decisions on twins: severity-scaled regret}\label{sec:regret}
Keeping a step open because closure cannot be certified may be wrong in a world where the
mechanism is in fact fault-free. When evidence-indistinguishable worlds prefer opposite actions,
evidence cannot remove this dilemma. Unverifiability alone does not create it: if indistinguishable
mechanisms of fault mass $0.8$ and $0.9$ face an open path of escape $0.1$, always keeping the step
open has zero regret.

\begin{theorem}[Twin regret]\label{thm:regret}
Consider two worlds whose evidence records have laws $Q_0,Q_1$ with $\TV(Q_0,Q_1)\le\varepsilon$, and
suppose that in each world the deployment losses are independent of the evidence record. In world~0,
closure beats the open path by expected total loss $\alpha S>0$; in world~1, the open path beats
closure by $\beta S>0$. Every decision rule based on the evidence, possibly randomized, has
\[
\max\{R_0,R_1\}\ \ge\ (1-\varepsilon)\,\frac{\alpha\beta}{\alpha+\beta}\,S,
\]
where $R_j$ is the expected regret in world $j$, measured against the better action's expected loss. For exact twins, $\varepsilon=0$ and the bound is
attained by choosing closure with probability $\alpha/(\alpha+\beta)$.
\end{theorem}
\begin{proof}
Let $p_j$ be the probability of choosing closure in world $j$. The data-processing property of total
variation gives $|p_0-p_1|\le\varepsilon$. Because losses are independent of the evidence, the expected
loss of a rule is $p_j$ times the closed loss plus $1-p_j$ times the open loss, so the regrets are $R_0=\alpha S(1-p_0)$ and $R_1=\beta Sp_1$,
so $R_0/(\alpha S)+R_1/(\beta S)=1-p_0+p_1\ge1-\varepsilon$. Since
$\max\{R_0,R_1\}(1/(\alpha S)+1/(\beta S))\ge R_0/(\alpha S)+R_1/(\beta S)$, the bound follows. With
$\varepsilon=0$ and $p_0=p_1=\alpha/(\alpha+\beta)$, both regrets equal $\alpha\beta S/(\alpha+\beta)$.
\end{proof}
The theorem combines two classical results. The inequality $1-p_0+p_1\ge1-\varepsilon$ is Le~Cam's
two-point argument~\cite{lecam73}. The exact-twin rule and its regret are the two-state case of
minimax-regret treatment allocation~\cite{manski05,manski09}, with $\alpha S$ and $\beta S$ as the
welfare gaps. What is added here is their reading as the cost of the second arm.
The independence assumption matters: if the evidence predicts deployment losses, a rule can exploit
it. With a fair evidence bit $X$ and losses (closed, open) equal to $(0,2)$ or $(1,0)$ in world~0 and
$(0,1)$ or $(2,0)$ in world~1 according to $X$, the unconditional preferences reverse, yet closing
exactly when $X=0$ has zero loss in both worlds. Under the assumption, the unavoidable worst-case regret
grows linearly with severity. A richer scheduler cannot resolve
it with evidence that remains indistinguishable; when the transcript includes an adaptive
build-and-test process, the twin relation is taken over that whole transcript. Only evidence that
separates the worlds, or a change of available actions such as an exogenous gate, reduces it. This
is the precise cost behind the second arm of the principle.

\section{Expected loss: thresholds for both arms}\label{sec:arms}
Admissibility uses a hard budget. A softer formulation compares expected severity-weighted loss plus
cost. Let $C_{\mathrm{op}}$ and $C_{\mathrm{cl}}$ be the expected costs of the open and closed paths
over the $n$ units, including build, validation, upkeep, gates, retries, and fallback.

\begin{theorem}[Two arms under expected loss]\label{thm:arms}
Let the open path be in $\mathcal P(\ell,u)$ and let $\varphi$ be a loss function.
\begin{enumerate}[label=(\alph*),itemsep=2pt]
\item \emph{Mandatory closure (ex ante).} Under the assumptions of
Theorem~\ref{thm:admissible}(iv), put
$\Delta_1=\E\varphi(\Bin(n,\ell))-(1-\eta)\E\varphi(\Bin(n,U))-\eta\varphi(n)$. If $\Delta_1>0$,
closure has lower ex-ante expected loss plus cost than the open path whenever
$S>S_1^\ast=(C_{\mathrm{cl}}-C_{\mathrm{op}})_+/\Delta_1$. If instead the uniform premise of
Theorem~\ref{thm:admissible}(iii) holds, take $\eta=0$; the comparison then holds for every compatible
model, not only in expectation. A statistical certificate with $\eta=0$ alone does not give this.
\item \emph{Forbidden closure (worst case).} Under the assumptions of
Theorem~\ref{thm:admissible}(v), put $\Delta_2=\E\varphi(\Bin(n,\tau(M)))-\E\varphi(\Bin(n,u))$. If
$\Delta_2>0$, the worst-case expected loss plus cost of closure over the compatible twins exceeds that
of the open path whenever $S>S_2^\ast=(C_{\mathrm{op}}-C_{\mathrm{cl}})_+/\Delta_2$.
\end{enumerate}
Both thresholds depend on the executors only through $\ell$ and $u$. A threshold exists only when
the corresponding gap is positive: a valid but loose certificate with $U\ge\ell$ gives no mandate.
\end{theorem}
\begin{proof}
(a) The open path's expected loss is at least $S\E\varphi(\Bin(n,\ell))$ by
Theorem~\ref{thm:dominance}, and the closed path's is at most
$S[(1-\eta)\E\varphi(\Bin(n,U))+\eta\varphi(n)]$ by Theorem~\ref{thm:admissible}(iv). Compare the
totals. (b) The worst case over compatible twins of the closed loss is at least $S\E\varphi(\Bin(n,\tau(M)))$,
and the open loss is at most $S\E\varphi(\Bin(n,u))$. Compare the totals.
\end{proof}
Part (b) is a worst-case preference. It does not say that the open path is better in every
compatible world: with twins of fault mass $0$ and $0.9$, an open path of escape $0.1$, equal costs,
and linear loss, (b) prefers the open path while the fault-free twin favours closure.
Theorem~\ref{thm:regret} prices that gap. For a fully determined step, $U=0$ and $\eta=0$, so
$\Delta_1=\E\varphi(\Bin(n,\ell))>0$ whenever $\ell>0$ and $\varphi(1)>0$, and $S_1^\ast$ is finite
for every cost difference.

\paragraph{Mean-only certificates.}
Validation may certify average fault mass without certifying a conditional ceiling.

\begin{proposition}[Sharp convex bounds from the mean]\label{prop:meanonly}
Let $N\in\{0,\dots,n\}$ with $\E N=m<n$, and let $\varphi$ be convex. With $k=\lfloor m\rfloor$ and
$\alpha=m-k$,
\[
(1-\alpha)\varphi(k)+\alpha\varphi(k+1)\ \le\ \E\varphi(N)\ \le\ \Big(1-\frac mn\Big)\varphi(0)+\frac mn\varphi(n),
\]
and both bounds are attained. If $\varphi$ is a nondecreasing convex loss and $\E N\le nU$, then
$\E\varphi(N)\le U\varphi(n)$, and closure has lower expected loss plus cost than an open path in
$\mathcal P(\ell,u)$ whenever $S[\E\varphi(\Bin(n,\ell))-U\varphi(n)]>C_{\mathrm{cl}}-C_{\mathrm{op}}$.
\end{proposition}
\begin{proof}
The lower bound is Jensen's inequality for the piecewise-linear interpolation of $\varphi$, attained
by a law on $\{k,k+1\}$. The upper bound is the chord of $\varphi$ between $0$ and $n$, attained by a
law on $\{0,n\}$. Each extremal law is the count of an exchangeable Bernoulli sequence: choose a
uniformly random subset of the drawn size. For the last claim, $\varphi(0)=0$ and the chord bound is
nondecreasing in $m$. Both bounds are classical: the upper one is the Edmundson--Madansky
bound~\cite{madansky59}, and both are the discrete convex extrema with a given mean~\cite{denuit97}.
\end{proof}
The gap between $\E\varphi(\Bin(n,U))$ and $U\varphi(n)$ measures what is lost when validation
cannot certify how faults cluster.

\section{Persistence and convex order}\label{sec:convex}
Linear loss sees only the mean escape. A convex loss also sees how escapes cluster. Part~V showed
that persistence, a shared defect or a latent regime, changes the damage distribution at equal
mean~\cite{p5}. Convex order makes this exact. For integer random variables write $X\cx Y$ if
$\E g(X)\le\E g(Y)$ for every convex $g$; this forces equal means. For a function $f$ on
$\{0,\dots,n\}$, the Bernstein polynomial is
$B_nf(\theta)=\sum_{k=0}^nf(k)\binom nk\theta^k(1-\theta)^{n-k}=\E f(\Bin(n,\theta))$.

\begin{lemma}[Bernstein polynomials preserve convexity]\label{lem:bernstein}
For $n\ge2$,
\[
(B_nf)''(\theta)=n(n-1)\sum_{k=0}^{n-2}\Delta^2f(k)\binom{n-2}{k}\theta^k(1-\theta)^{n-2-k},
\qquad\Delta^2f(k)=f(k+2)-2f(k+1)+f(k).
\]
Hence $B_nf$ is convex on $[0,1]$ if $f$ is convex, and strictly convex if every $\Delta^2f(k)>0$.
\end{lemma}
\begin{proof}
Differentiate twice and collect terms; this is the classical identity behind convexity
preservation~\cite{lorentz}.
\end{proof}

\begin{theorem}[Persistence is convex spread]\label{thm:convex}
Let $J_1,\dots,J_n$ be conditionally independent Bernoulli$(\Theta)$ given a random
$\Theta\in[0,1]$, and let $J'$ be the same with $\Theta'$. If $\Theta\cx\Theta'$, then $N_n\cx N'_n$.
At a given mean escape $\bar\theta$:
\begin{enumerate}[label=(\alph*),itemsep=2pt]
\item fresh independent errors, $\Theta\equiv\bar\theta$, are convex-minimal among such processes;
\item a closed mechanism drawn from a build process, deployed ungated with independent inputs so that
all units share $\Theta=\pi(M)$, has $\Bin(n,\bar\theta)\cx N_n$, strictly for strictly convex losses
when $n\ge2$ and $\pi(M)$ is not almost surely constant;
\item linear loss cannot distinguish them: $\E N_n=n\bar\theta$ in both cases, while
$\operatorname{Var}N_n=n\bar\theta(1-\bar\theta)+n(n-1)\operatorname{Var}\Theta$.
\end{enumerate}
\end{theorem}
\begin{proof}
For convex $g$, $\E g(N_n)=\E[B_ng(\Theta)]$, and $B_ng$ is convex by Lemma~\ref{lem:bernstein}, so
$\Theta\cx\Theta'$ gives $\E g(N_n)\le\E g(N'_n)$. Part (a): a constant is convex-below every variable
with the same mean, by Jensen. Part (b) is (a) with strict Jensen for the strictly convex $B_ng$. Part
(c): $B_n$ maps linear functions to linear functions, and the variance follows from the law of total
variance.
\end{proof}
The implication from $\Theta\cx\Theta'$ to $N_n\cx N'_n$ for mixed binomial counts is
known~\cite{adell,lefevre17}; the new part is its reading as persistence against fresh errors.
Part~(a) concerns mixtures only: independent errors with unequal probabilities and the same mean are
convex-below $\Bin(n,\bar\theta)$~\cite{hoeffding56}.

\begin{remark}
Three limits apply. First, the comparison is within conditionally independent processes; finite
exchangeability is not enough: a uniformly random ordering of $(0,1)$ is exchangeable with constant
count $1$, convex-below $\Bin(2,1/2)$. Second, the mixture over builds expresses uncertainty about
which build is deployed; for a single deployed build the count is binomial with its own fault mass.
A persistent executor regime drawn once per run (Part~V) is also a mixture, so the contrast is
persistent versus fresh, not mechanism versus executor. Third, the any-escape indicator is not
convex and moves the other way: with $n=2$ and mean $1/2$, complete persistence gives
$\Prb(N\ge1)=1/2$ against $3/4$ for fresh errors, while $\Prb(N=2)$ rises from $1/4$ to $1/2$.
Persistence trades frequent small damage for rare concentrated damage; the loss shape decides which
matters.
\end{remark}

\begin{example}[Same mean, different risk]\label{ex:spread}
Take $n=10$, mean escape $0.1$, and $\varphi(k)=\max(0,k-2)^2$, a loss that starts after two escapes.
Fresh errors give $\E\varphi=0.1179$. A mechanism with fault mass $0.02$ or $0.18$ with equal
probability gives $0.3327$. Under $\varphi(k)=k$ both give $1.0000$; the variances are $0.900$ and
$1.476$.
\end{example}

\section{An anytime budget for a scheduler}\label{sec:budget}
A scheduler chooses actions unit by unit, and the number of units may depend on outcomes. Local
certificates then need to add up to a lifetime guarantee.

\begin{theorem}[Adaptive consequence budget]\label{thm:budget}
At each unit the scheduler chooses, from the history, an action, a weight $s_i\ge0$, and a
conditional escape ceiling $u_i$ with $\Prb(J_i=1\mid\mathcal F_{i-1})\le u_i$. Let
$D_t=\sum_{i\le t}s_iJ_i$ and let $T$ be a stopping time.
\begin{enumerate}[label=(\alph*),itemsep=2pt]
\item If $\sum_{i\le T}s_iu_i\le B$ almost surely, then $\E D_T\le B$. If $s_i$ bounds the conditional
expected damage of an escape at unit $i$ given the history, the expected total damage is at most $B$.
\item For every fixed $\lambda>0$ and $\delta\in(0,1)$,
\[
\Prb\Big(\exists t:\ \lambda D_t-\sum_{i\le t}\log\big[1+u_i(e^{\lambda s_i}-1)\big]\ge\log(1/\delta)\Big)\le\delta .
\]
\end{enumerate}
In (b), $D_t$ bounds the actual damage pathwise only if each $s_i$ bounds the damage of unit $i$
pathwise. Both parts hold under adaptive routing and every mixture of executors whose ceilings hold on
the routed histories.
\end{theorem}
\begin{proof}
(a) $\E D_T=\E\sum_{i\le T}s_i\Prb(J_i=1\mid\mathcal F_{i-1})\le\E\sum_{i\le T}s_iu_i\le B$, by
optional stopping for the nonnegative increments with predictable weights; monotone convergence
handles unbounded $T$. For the damage claim, the expected damage of unit $i$ given the history is
at most $\Prb(J_i=1\mid\mathcal F_{i-1})s_i$. (b) $\E[e^{\lambda s_iJ_i}\mid\mathcal F_{i-1}]=1+\Prb(J_i=1\mid\mathcal
F_{i-1})(e^{\lambda s_i}-1)\le1+u_i(e^{\lambda s_i}-1)$, so
$M_t=\exp(\lambda D_t-\sum_{i\le t}\log[1+u_i(e^{\lambda s_i}-1)])$ is a nonnegative supermartingale
with $M_0=1$. Ville's maximal inequality gives the bound~\cite{ville}. This is the standard
supermartingale construction of time-uniform Chernoff bounds~\cite{howard20}.
\end{proof}
For expected damage, $s_i$ can be the structural severity of Section~\ref{sec:severity}, an expected
damage. For the anytime tail of actual damage, $s_i$ must be a pathwise bound, such as the total cost of
the reach set; with expected severities, part (b) controls the severity-weighted escape sum, a proxy. The theorem turns per-step certificates into a budget that the
scheduler can spend: closing a step with a verified mechanism lowers its $u_i$; adding a gate lowers
$u_i$ at a cost; high-severity steps consume the budget fastest.

\section{Producers average, reviewers intersect}\label{sec:review}
\begin{proposition}[Production mixing and review]\label{prop:review}
Let producers $e$ have badness $p_e$ and fault laws $\mu_e$, and let gates accept their blind sets at
once.
\begin{enumerate}[label=(\alph*),itemsep=2pt]
\item \begin{sloppypar}Routing production with weights $w_e$ gives the first-attempt floor
$\sum_ew_ep_e\mu_e(\Sigma)\ge\min_ep_e\mu_e(\Sigma)$; lower bounds average as lower bounds.\end{sloppypar}
\item Conjunctive review by $r_1,\dots,r_k$ gives producer $e$ the escape
$p_e\mu_e(\bigcap_j\Sigma_{r_j})$ if the blind sets are complete, and at least that if they are only
guaranteed.
\item \emph{Universal superiority.} Let costs be equal, let blind sets be complete, and for every fault
$z$ let some admissible producer have $p_e>0$ and $\mu_e=\delta_z$. Reviewer $r_1$ has escape at most that of $r_2$ for every
producer if and only if $\Sigma_{r_1}\subseteq\Sigma_{r_2}$. More generally, let a unit's state $z$ range
over a fault-free state $z_0$ and the fault states, with residual damages $h_j\ge0$, $h_j(z_0)=0$, and
per-unit costs $c_j$, so that design $j$'s per-unit total is $c_j+S\,\E h_j(Z)$ (for producer $e$,
$\E h_j(Z)=p_e\int h_j\,d\mu_e$). If every point-mass state law is admissible, including the fault-free
state, design~1 is at least as good for every admissible state law if and only if
$c_1+Sh_1(z)\le c_2+Sh_2(z)$ for every state $z$, which requires $c_1\le c_2$. With a smaller admissible
class only sufficiency holds: if every producer has $p_e=0.01$, then with $S=1$, $(c_1,h_1)=(0.005,0)$
and $(c_2,h_2)=(0,\ind)$, design~1 costs $0.005<0.01$ for every such producer although it is worse at
the fault-free state.
\item \emph{No universal ranking.} Let blind sets be complete, and let producers with positive badness
concentrated on $\Sigma_{r_1}\setminus\Sigma_{r_2}$, and others concentrated on
$\Sigma_{r_2}\setminus\Sigma_{r_1}$, be admissible. Then each reviewer is strictly better for some producer.
Both conditions are needed: with guaranteed sets $\{x\}$ and $\{y\}$ for reviewers that in fact accept
both faults, the escapes are equal for every producer; and if every admissible producer has fault law
$(0.9,0.1)$ on $(x,y)$, the reviewer missing only $y$ is always better.
\item \emph{Coalitions.} Pairwise single-reviewer escapes do not determine a coalition's escape: two
reviewers each missing half of a producer's faults can have an intersection of mass $0$ or $1/2$.
\end{enumerate}
\end{proposition}
\begin{proof}
(a) The floor is linear in the fault law. (b) A fault passes only if every reviewer accepts it.
(c) Integrate the pointwise inequality for sufficiency; for necessity, take a point mass at a state
that violates it. The reviewer case has $h_j=\ind_{\Sigma_{r_j}}$ and equal costs. Omitting the fault-free
state is an error: with $p_e=0.01$, $S=2$, $(c_1,h_1)=(1,0)$ and $(c_2,h_2)=(0,\ind)$, the fault states
alone favour design~1, while the per-unit totals are $1$ and $0.02$. (d) Take an admissible producer concentrated on $\Sigma_{r_2}\setminus\Sigma_{r_1}$: with complete sets,
$r_1$ rejects its faults and $r_2$ accepts them; symmetrically for the other difference. (e) Take faults uniform on
$\{1,2,3,4\}$, $\Sigma_{r_1}=\{1,2\}$, and $\Sigma_{r_2}=\{3,4\}$ or $\{1,2\}$.
\end{proof}
Producer floors are governed by bounds that survive every mixture. Reviewer rankings are not
invariant to the producer's error law, except under blind-set inclusion. A scheduler therefore needs
review data indexed by producer and by reviewer set, not a single quality score per reviewer and not
only pairwise data. Parts (b) and (e) are the gated form of design-diversity models of software
reliability: the probability that diverse versions fail together is the demand-profile mass of the
intersection of their failure sets, and single-version failure probabilities do not determine
it~\cite{eckhardt85,littlewood89,lps01}. Part (c) is statewise dominance: an option has lower
expected loss under every admissible law if and only if its loss is pointwise lower. A reviewer from
the producer's own family inherits the producer's blind set (Part~II). Errors of language models are
substantially correlated in practice, partly through shared architectures and providers, and for
larger models also across providers~\cite{kim25}. Part (c) also shows that a gate can be preferable because it blocks the most damaging
faults, even if it misses more faults by count.

\section{Severity from structure}\label{sec:severity}
\begin{definition}[Structural severity]\label{def:severity}
Let the work form a directed acyclic dependency graph. Gates may sit on edges. A fault that escapes
at node $v$ reaches node $w$ with probability $\rho(v\to w)$, with $\rho(v\to v)=1$. Reaching $w$ costs
$c(w)\ge0$ unless the damage is undone; let $\iota(w)\in[0,1]$ be the probability that damage at $w$ is
not undone, given that $w$ is reached, and assume this event is independent of how $w$ was reached. The
severity of $v$ is the expected damage of one escape at $v$,
\[
S(v)=\sum_{w}\rho(v\to w)\,c(w)\,\iota(w).
\]
It is an expectation, not a bound on realized damage: if an escape reaches a node of cost $100$ with
probability $0.01$, then $S=1$ while the realized damage can be $100$.
\end{definition}
The reach set $\{w:\rho(v\to w)>0\}$ is the blast radius of $v$; $S$ weights it by cost and
irreversibility. With $c\iota\equiv1$, $S(v)$ is the expected spread of an influence-propagation
model~\cite{kempe03}; weighting reach by consequence follows the quantitative definition of risk as
probability and consequence~\cite{kaplan81}.

\begin{proposition}[Properties of structural severity]\label{prop:severity}
\begin{enumerate}[label=(\alph*),itemsep=2pt]
\item \emph{Factorization and its limits.} In general, the expected damage of producer $e$ at $v$ is
$R_e(v)=\int h(v,z)\,\nu_e(dz)$, where $\nu_e$ is the escaped-fault subprobability measure and
$h(v,z)$ the downstream damage of a fault of type $z$. It factorizes as $\nu_e(\text{all})\,S(v)$ with
an executor-independent $S(v)$ exactly when the conditional expected damage of an escape is the same
for every executor, for example when propagation, interception, and repair do not depend on the
fault type. Without that condition, the \emph{damage envelope} $S_{\max}(v)=\sup_zh(v,z)$ gives
$R_e(v)\le\nu_e(\text{all})\,S_{\max}(v)$ for every executor and every mixture.
\item \emph{Ancestry bound.} If edges transmit faults independently and $\rho(v\to w)$ is the
probability that a transmitting path exists, then $S(v)\ge\rho(v\to w)S(w)$ for every descendant $w$.
If transmission is certain along every path, $S(v)\ge S(w)$. An exact gate between $v$ and $w$ makes
$\rho(v\to w)$ small, so faults from $v$ can be far less severe than faults born at $w$.
\item \emph{Additivity is an upper bound in expectation.} With several escapes, the sum of severities
bounds the expected damage from above when damage is the cost of the union of reached nodes, with
equality when the reached sets are disjoint almost surely.
\end{enumerate}
\end{proposition}
\begin{proof}
(a) The factorization is the definition of a conditional expectation that does not depend on $e$.
The envelope bound integrates $h\le S_{\max}$.
(b) For each node $x$, the events ``$v$ reaches $w$'' and ``$w$ reaches $x$'' are increasing in the
independent edge variables. Harris's inequality gives $\Prb(\text{both})\ge\rho(v\to w)\rho(w\to x)$
\cite{harris,fkg}, and ``$v$ reaches $x$'' contains their intersection. Multiply by $c(x)\iota(x)\ge0$
and sum over $x$. Certain transmission gives $\rho(v\to w)=1$.
(c) Pathwise, the cost of a union is at most the sum of costs, with equality for disjoint sets; take
expectations.
\end{proof}
Under the condition of part (a), risk factorizes into an executor-dependent escape law and a
structure-dependent severity, which a scheduler can compute once from the dependency graph and
reuse across executors. Otherwise the envelope gives an executor-invariant upper bound. Escape laws,
floors, and ceilings must still be estimated per executor.

\section{Relation to the earlier papers}\label{sec:relation}
\begin{proposition}[Reduction to the mechanization criterion]\label{prop:reduction}
Let $\varphi(k)=k$, let severity be a constant $S>0$, and let the expected per-unit escapes of the open
and closed paths be $o$ and $d$. Then closure has lower expected loss plus cost exactly when
$Sn(o-d)>C_{\mathrm{cl}}-C_{\mathrm{op}}$. With equal costs this is Part~IV's criterion $d<o$; with the
expected loss included in the per-call saving, it is Part~IV's break-even.
\end{proposition}
\begin{proof}
By linearity, expected losses are $Sno$ and $Snd$, without any independence assumption. Rearrange.
\end{proof}
Substituting bounds for $o$ and $d$ gives a sufficient criterion, not Part~IV's exact one. Severity
therefore changes the decision only through the cost trade-off, the confidence tail, the gaps between
bounds, and nonlinear loss. Theorem~\ref{thm:dominance} with the any-escape indicator recovers
Part~V's conditional floor. The invariance here is of the same kind as Part~V's: conditional risks,
not arrival laws, carry the argument.

\section{What must be measured}\label{sec:measure}
The theorems need four kinds of quantity, and bounds suffice for their decisions.
\begin{enumerate}[leftmargin=*,itemsep=2pt]
\item \emph{Executor floors and ceilings.} A floor lies in the blind set of the producing family's
gates, so those gates cannot measure it. An exogenous exact audit reveals escapes. A binomial
confidence bound from $n_a$ audited outputs, for example Clopper--Pearson~\cite{clopper}, estimates a
stationary escape parameter under its sampling model; it does not certify a uniform conditional
bound under arbitrary dependence or workload shift. Adaptive auditing and repeated looks need
time-uniform confidence sequences~\cite{howard}.
\item \emph{Certified mechanism levels.} With $m$ independent representative validation samples and
zero failures, the one-sided upper bound is $U=1-\eta^{1/m}$~\cite{p5}. This is a statistical bound,
not proof of zero fault mass, and selection among builds needs simultaneous inference. Evidence that
cannot separate twins cannot, with high probability, bring a fault-mass bound below $\tau_\pi$ or an
escape bound below $\tau$. A fault-mass certificate certifies closed-path escape when raw faults and
escapes are justifiably related: with independent target inputs, a gate that only suppresses outputs,
and no fallback, $J_i\le F_i$ for the raw-fault indicator $F_i$, so $\Prb(J_i=1\mid\mathcal F_{i-1})\le\pi(M)$
and a valid bound on $\pi(M)$ bounds escape with the same confidence. With a fallback or other
interaction, validation of the whole closed path, or a known structural bound, is needed.
\item \emph{Severity.} The dependency graph, costs, and irreversibility give $S(v)$ or its envelope.
\item \emph{Review coverage.} Escape must be measured per producer and per reviewer set
(Proposition~\ref{prop:review}).
\end{enumerate}

\section{Executable checks}\label{sec:checks}
A companion script (\texttt{simulate\_risk\_mechanization.py}) checks the finite claims: Bernstein convexity and its second-derivative identity,
binomial dominance by exact enumeration of history-dependent processes, mixture closure, convex order
of binomial mixtures, the mean-only bounds, the two-arm thresholds, the twin-certification and regret
bounds, the budget supermartingale, the reviewer and coalition examples, and the severity ancestry
bound by enumeration. Negative controls reinstate false claims and must be rejected.

\begin{center}\small
\begin{tabular}{p{0.66\textwidth}r}
\toprule Quantity & Value\\\midrule
Example~\ref{ex:spread}: fresh errors, $\E\varphi$ & 0.1179\\
Example~\ref{ex:spread}: mechanism mixture, $\E\varphi$ & 0.3327\\
Example~\ref{ex:spread}: linear loss, both & 1.0000\\
Example~\ref{ex:spread}: variance, fresh / mixture & 0.900 / 1.476\\
Theorem~\ref{thm:regret}: $\alpha=2$, $\beta=3$, $S=1$, $\varepsilon=0$: minimax regret & 1.2000\\
\bottomrule
\end{tabular}
\end{center}

\section{Related work and limitations}\label{sec:related}
Risk-based integrity requirements, in which consequence and likelihood set a required failure
measure, are standard in functional safety~\cite{iec61508}. Our budget is that idea in the
gated-computation setting. Choosing a level of automation from the reliability of the automation and
the cost of outcomes is an established design question~\cite{parasuraman00}, keeping an untrusted model
behind trusted monitoring is studied as AI control~\cite{greenblatt23}, and expected-loss deferral
thresholds go back to the reject option in classification~\cite{chow70}.

On the decision side, the worst case over compatible models is the maximum risk of Wald's minimax
criterion~\cite{wald45} over a set of models, as in multiple-prior decision theory and partial identification~\cite{gilboa89,manski00}.
The twin regret is Le~Cam's two-point bound~\cite{lecam73} combined with minimax-regret
allocation~\cite{manski05,manski09}. The limits of statistical validation for very high
dependability~\cite{littlewood93,butler93} are the practical content of the second arm. Reviewer
conjunction follows the models of coincident failure in diverse software~\cite{eckhardt85,littlewood89,lps01},
and the errors of language models are empirically correlated~\cite{kim25}. Structural severity is a
weighted reach, as in influence spread~\cite{kempe03}, with risk as probability and
consequence~\cite{kaplan81}.

On the probability side, domination by independent variables under conditional bounds is
classical~\cite{doerr}. Stochastic and convex orders and their preservation by Bernstein-type operators
are standard~\cite{shaked,adell}, convex order of mixed binomial counts is known~\cite{lefevre17},
heterogeneous independent counts are convex-below the binomial~\cite{hoeffding56}, and the sharp convex
bounds from the mean are the Edmundson--Madansky bound and its discrete
counterpart~\cite{madansky59,denuit97};
Bernstein convexity is classical~\cite{lorentz}; positive association goes back to
Harris~\cite{harris,fkg}; infinitely exchangeable sequences are mixtures of independent ones, and the
first $k$ terms of an exchangeable binary sequence of length $n$ are within total variation of order
$k/n$ of such a mixture, while a whole finite exchangeable vector can be far from every
mixture~\cite{diaconis};
Ville's inequality, supermartingale Chernoff bounds, and confidence sequences underlie the
budget~\cite{ville,howard20,howard}.

What we add is the following. Theorem~\ref{thm:admissible} states when closure is mandatory, when it
is forbidden, and when no option fits; we are not aware of a prior statement of this two-arm
structure. The classical bounds are read as conditions that hold uniformly over the executor's error
law, dependence, routing, and mixtures. Structural severity and its properties are stated for gated
work, and Proposition~\ref{prop:reduction} reduces the results to Part~IV's criterion.

The results compare options through bounds. They do not estimate floors, ceilings, or severity for
any real executor, and they do not optimize an adaptive policy. The convex-order results are within
conditionally independent processes. Severity factorizes only under the stated condition. The
mechanize and gate options interact with maintenance and drift, treated in Parts~IV and~V and not
here. This is a preprint. It has not been peer reviewed. The prior-art search was targeted, not
systematic.

\section{Conclusion}
Severity sets a budget, and the budget, compared with bounds that hold for every compatible error law,
decides. Closure is mandatory when it is admissible and every other available option is not; with the
budget below every executor floor, an open path with inherited gates is excluded, and only a closure
that fits the budget or a gate outside the family remains. Closure is forbidden when an evidence-compatible closed model violates the budget, and when
indistinguishable worlds prefer opposite actions, every evidence-based rule pays a worst-case regret
proportional to severity. When no option fits, the step needs distinguishing evidence, an exogenous
gate, or a redesign. These conditions use only floors, ceilings, certified levels, and compatible twins,
so they hold uniformly over error laws, routing, and mixtures of executors within the stated bounds.
Persistence matters beyond the mean through convex order, review value is producer-relative unless
blind sets are nested, and a martingale budget turns local certificates into a scheduler's lifetime
guarantee in expectation, and in tail when damage bounds are pathwise.

\section*{Code and data availability}
The standard-library Python script \texttt{simulate\_risk\_mechanization.py} reproduces the finite checks and the
values of Section~\ref{sec:checks}. It is at \url{https://github.com/ivmat/gated-computation-sim}.
This paper is archived on Zenodo, DOI~\href{https://doi.org/10.5281/zenodo.23324667}{10.5281/zenodo.23324667}.

\begin{thebibliography}{99}
\bibitem{p1} I. Matija\v{s}evi\'c. A Fault-Tolerance Threshold for Gated Agentic Computation.
2026. \url{https://doi.org/10.5281/zenodo.20820968}.
\bibitem{p2} I. Matija\v{s}evi\'c. Generated Gates Inherit Their Generator's Blind Spots.
2026. \url{https://doi.org/10.5281/zenodo.20837102}.
\bibitem{p3} I. Matija\v{s}evi\'c. Escape, Cost, and Correlation at the Verification Floor of
Gated Agentic Computation. 2026. \url{https://doi.org/10.5281/zenodo.21456783}.
\bibitem{p4} I. Matija\v{s}evi\'c. Progressive Semantic Mechanization. 2026.
\url{https://doi.org/10.5281/zenodo.23119225}.
\bibitem{p5} I. Matija\v{s}evi\'c. Beyond Constant Fault Rates: Reliability Bounds for Gated Computation and Semantic Mechanization. 
Preprint, 2026. \url{https://doi.org/10.5281/zenodo.23293740}.
\bibitem{lorentz} G. G. Lorentz. \emph{Bernstein Polynomials}. University of Toronto Press, 1953.
\bibitem{adell} J. A. Adell and A. P\'erez-Palomares. Stochastic orders in preservation properties by
Bernstein-type operators. \emph{Advances in Applied Probability} 31(2):492--509, 1999.
\bibitem{harris} T. E. Harris. A lower bound for the critical probability in a certain percolation
process. \emph{Proceedings of the Cambridge Philosophical Society} 56:13--20, 1960.
\bibitem{fkg} C. M. Fortuin, P. W. Kasteleyn, and J. Ginibre. Correlation inequalities on some
partially ordered sets. \emph{Communications in Mathematical Physics} 22:89--103, 1971.
\bibitem{shaked} M. Shaked and J. G. Shanthikumar. \emph{Stochastic Orders}. Springer, 2007.
\bibitem{diaconis} P. Diaconis and D. Freedman. Finite exchangeable sequences.
\emph{Annals of Probability} 8(4):745--764, 1980.
\bibitem{clopper} C. J. Clopper and E. S. Pearson. The use of confidence or fiducial limits
illustrated in the case of the binomial. \emph{Biometrika} 26(4):404--413, 1934.
\bibitem{ville} J. Ville. \emph{\'Etude critique de la notion de collectif}. Gauthier-Villars, 1939.
\bibitem{howard} S. R. Howard, A. Ramdas, J. McAuliffe, and J. Sekhon. Time-uniform, nonparametric,
nonasymptotic confidence sequences. \emph{Annals of Statistics} 49(2):1055--1080, 2021.
\bibitem{iec61508} IEC 61508. Functional safety of electrical/electronic/programmable electronic
safety-related systems. International Electrotechnical Commission, 2010.
\bibitem{doerr} B. Doerr. Probabilistic tools for the analysis of randomized optimization heuristics.
In B. Doerr and F. Neumann (eds.), \emph{Theory of Evolutionary Computation}, Springer, 2020,
pp.~1--87. \url{https://doi.org/10.1007/978-3-030-29414-4_1}.
\bibitem{wald45} A. Wald. Statistical decision functions which minimize the maximum risk.
\emph{Annals of Mathematics} 46(2), 1945. \url{https://doi.org/10.2307/1969022}.
\bibitem{gilboa89} I. Gilboa and D. Schmeidler. Maxmin expected utility with non-unique prior.
\emph{Journal of Mathematical Economics} 18(2):141--153, 1989.
\bibitem{manski00} C. F. Manski. Identification problems and decisions under ambiguity: empirical
analysis of treatment response and normative analysis of treatment choice. \emph{Journal of
Econometrics} 95(2):415--442, 2000.
\bibitem{lecam73} L. Le~Cam. Convergence of estimates under dimensionality restrictions. \emph{Annals
of Statistics} 1(1), 1973. \url{https://doi.org/10.1214/aos/1193342380}.
\bibitem{manski05} C. F. Manski. Fractional treatment rules for social diversification of indivisible
private risks. NBER Working Paper 11675, 2005.
\bibitem{manski09} C. F. Manski. The 2009 Lawrence R. Klein lecture: diversified treatment under
ambiguity. \emph{International Economic Review} 50(4):1013--1041, 2009.
\bibitem{littlewood93} B. Littlewood and L. Strigini. Validation of ultrahigh dependability for
software-based systems. \emph{Communications of the ACM} 36(11):69--80, 1993.
\bibitem{butler93} R. W. Butler and G. B. Finelli. The infeasibility of quantifying the reliability of
life-critical real-time software. \emph{IEEE Transactions on Software Engineering} 19(1):3--12, 1993.
\bibitem{madansky59} A. Madansky. Bounds on the expectation of a convex function of a multivariate
random variable. \emph{Annals of Mathematical Statistics} 30(3):743--746, 1959.
\bibitem{denuit97} M. Denuit and C. Lef\`evre. Some new classes of stochastic order relations among
arithmetic random variables, with applications in actuarial sciences. \emph{Insurance: Mathematics and
Economics} 20(3):197--213, 1997.
\bibitem{lefevre17} C. Lef\`evre, S. Loisel, and S. Utev. On finite exchangeable sequences and their
dependence. \emph{Journal of Multivariate Analysis} 162:93--109, 2017.
\bibitem{hoeffding56} W. Hoeffding. On the distribution of the number of successes in independent
trials. \emph{Annals of Mathematical Statistics} 27(3):713--721, 1956.
\bibitem{howard20} S. R. Howard, A. Ramdas, J. McAuliffe, and J. Sekhon. Time-uniform Chernoff bounds
via nonnegative supermartingales. \emph{Probability Surveys} 17:257--317, 2020.
\bibitem{eckhardt85} D. E. Eckhardt and L. D. Lee. A theoretical basis for the analysis of multiversion
software subject to coincident errors. \emph{IEEE Transactions on Software Engineering}
SE-11(12):1511--1517, 1985.
\bibitem{littlewood89} B. Littlewood and D. R. Miller. Conceptual modeling of coincident failures in
multiversion software. \emph{IEEE Transactions on Software Engineering} 15(12):1596--1614, 1989.
\bibitem{lps01} B. Littlewood, P. Popov, and L. Strigini. Modeling software design diversity: a
review. \emph{ACM Computing Surveys} 33(2):177--208, 2001.
\bibitem{kim25} E. Kim, A. Garg, K. Peng, and N. Garg. Correlated errors in large language models.
2025. \url{https://arxiv.org/abs/2506.07962}.
\bibitem{kempe03} D. Kempe, J. Kleinberg, and \'E. Tardos. Maximizing the spread of influence through a
social network. In \emph{Proceedings of KDD}, pp.~137--146, 2003.
\bibitem{kaplan81} S. Kaplan and B. J. Garrick. On the quantitative definition of risk. \emph{Risk
Analysis} 1(1):11--27, 1981.
\bibitem{parasuraman00} R. Parasuraman, T. B. Sheridan, and C. D. Wickens. A model for types and levels
of human interaction with automation. \emph{IEEE Transactions on Systems, Man, and Cybernetics,
Part~A} 30(3):286--297, 2000.
\bibitem{greenblatt23} R. Greenblatt, B. Shlegeris, K. Sachan, and F. Roger. AI control: improving
safety despite intentional subversion. 2023. \url{https://arxiv.org/abs/2312.06942}.
\bibitem{chow70} C. K. Chow. On optimum recognition error and reject tradeoff. \emph{IEEE Transactions
on Information Theory} 16(1):41--46, 1970.
\end{thebibliography}
\end{document}
